% ============================================================
% f2015ahl_q13_flowchart.tex —— 2015 安徽卷（理）第 13 题 程序框图
%
% 【来源】工作档 .work/phase3-b13/source/2015安徽理.pdf（sha256 402eca00…ce7b0）
%          第 1 页（页脚 931）栏 2；坐标数据取自实测清单
%          .work/phase3-b13/figure-inventory.md §F03（要素表 P094–P115、标签表）。
% 【坐标口径】真页 pt（PDF 页面坐标，原点页左下角，1190.55 × 841.89 bp）。
%          画布 1 单位 = 1 TeX pt，图内坐标直接写字面值（清单已给真页 pt）。
%          清单 §0.3 本卷实测换算因子 a_x = 0.999622795、a_y = 0.999621571，
%          与 1 的偏差 3.8e-4 系渲染器取整所致（非页面缩放），故不再二次缩放。
%          注：1 TeX pt = 0.996264 bp，故本图绝对尺寸比源件小 0.375%（项目定型口径）。
% 【线宽】框体、连接线、折线一律 0.3980 pt；框内分式横线 0.4200 pt（源件两档混用，
%          清单 FI-S10）——分式横线由 `\frac` 内生，不另画路径。
% 【圆角】两个终端框圆角半径 3.985 pt，源件四角为 κ = 0.552321 的贝塞尔四分之一圆弧
%          （标准 κ = 0.5522847，偏差 0.0066 %），用 `rounded corners` 复刻。
% 【照源复刻的存疑项】
%   1) 回流折线并入【主脊竖线】而非汇入菱形左顶点（清单 FI-S8）：折线末段止于
%      x = 578.6640（缺口 3.587 pt 由填充箭头本体补齐），箭头尖端 (582.2510,289.3726)
%      落在主脊竖线 P104（y∈[285.8723,296.4593]）上，高出菱形顶顶点 7.37 pt —— 照源保留。
%   2) 菱形左顶点 (508.3574,263.5927) 处无任何回边端点，照源保留。
%   3) 判据文字为十进制 `1.414`（清单 §F03.2(f)）与 `0.005`，照源排版。
%   4) 三个矩形框源件用 `re` 算子、菱形/平行四边形/圆角框用 m+l+h（FI-S9），
%      属绘制路径差异、几何无差，此处统一用 TikZ 图元。
%   5) 标签按【源墨迹实测中心】放置；清单给的是 pdftohtml em 框中心（含 descend，
%      中心偏低约 1.0 pt），故坐标 = 清单值 + 墨迹修正。
% ============================================================
% 【自验】（探针 .work/phase3-b13/probe/q13-probe.tex，latexmk -pdfxe）
%   编译：0 error、0 Overfull、0 Underfull、0 Missing character。
%   字面量反算双层核对（A 手工誊录清单字面量 + B 源 PDF 逐命令复解析，含两条
%   曲线的全部 598 段控制点）：最大残差 0.0000 pt（阈值 0.05）。
%   600 dpi 墨迹核对（源件页 vs 探针对位后）：几何 213 点采样 max 0.245 pt / rms 0.077 pt，标签墨迹质心 max 0.093 pt；
%   并排与红蓝重合裁片见 probe/crops/q13-side.png、q13-overlay.png。
\begin{tikzpicture}[x=1pt, y=1pt,
    every node/.style={font=\fontsize{10.5pt}{12.6pt}\selectfont, inner sep=0pt}]
  \begin{scope}[line width=0.398pt]
    % ---------- 框体（§F03.2(a)）----------
    \draw[rounded corners=3.985pt] (568.2517,327.0230) rectangle (596.2507,343.8610);  % P094 开始（圆角）
    \draw (550.3670,296.6590) rectangle (614.1350,312.4520);                           % P095 a=1,n=1
    \draw (582.2514,282.0037) -- (508.3574,263.5927) -- (582.2514,245.1817)            % P096 菱形
          -- (656.1444,263.5927) -- cycle;
    \draw (546.9220,201.5800) rectangle (617.5800,230.5270);                           % P097 a=1+1/(1+a)
    \draw (556.4070,172.3670) rectangle (608.0950,187.0080);                           % P099 n=n+1
    \draw (643.2703,207.6238) -- (649.4063,224.4828) -- (693.3403,224.4828)            % P100 输出 n（平行四边）
          -- (687.2043,207.6238) -- cycle;
    \draw[rounded corners=3.985pt] (654.3089,176.0349) rectangle (682.3079,193.0509);  % P101 结束（圆角）
    % ---------- 连接线与折线（§F03.2(c)）----------
    \draw (582.2510,326.8238) -- (582.2510,316.2368);   % P102 开始 ↓ 初值框
    \draw (582.2510,296.4593) -- (582.2510,285.8723);   % P104 初值框 ↓ 菱形
    \draw (582.2510,244.9002) -- (582.2510,234.3122);   % P106 菱形 ↓ 赋值框
    \draw (582.2510,201.3805) -- (582.2510,190.7935);   % P108 赋值框 ↓ 自增框
    \draw (668.3081,207.4239) -- (668.3071,196.8369);   % P110 输出框 ↓ 结束框
    \draw (582.2510,172.1676) -- (582.2510,160.8286) -- (497.2110,160.8286)             % P112 回流折线（4 段）
          -- (497.2110,289.3726) -- (578.6640,289.3726);
    \draw (656.4261,263.5927) -- (668.3071,263.5927) -- (668.3071,228.2707);            % P114 「否」支（2 段）
    % ---------- 7 个填充箭头（尖端坐标取清单；Triangle 长 3.985 / 宽 2.988）----------
    \draw[{Triangle[length=3.985pt,width=2.988pt]}-] (582.2510,312.6504) -- (582.2510,316.6354);  % P103 ↓
    \draw[{Triangle[length=3.985pt,width=2.988pt]}-] (582.2510,282.2861) -- (582.2510,286.2711);  % P105 ↓
    \draw[{Triangle[length=3.985pt,width=2.988pt]}-] (582.2510,230.7261) -- (582.2510,234.7111);  % P107 ↓
    \draw[{Triangle[length=3.985pt,width=2.988pt]}-] (582.2510,187.2072) -- (582.2510,191.1922);  % P109 ↓
    \draw[{Triangle[length=3.985pt,width=2.988pt]}-] (668.3074,193.2502) -- (668.3075,197.2352);  % P111 ↓
    \draw[{Triangle[length=3.985pt,width=2.988pt]}-] (582.2510,289.3726) -- (578.2660,289.3726);  % P113 → 回边末端
    \draw[{Triangle[length=3.985pt,width=2.988pt]}-] (668.3075,224.6844) -- (668.3075,228.6694);  % P115 ↓
  \end{scope}
  % ---------- 文字标签（源墨迹实测中心；行末为清单值，供追溯）----------
  \node at (582.096,335.499) {开始};                       % L q13_start  核字框中心 (582.936,335.555)
  \node at (582.209,304.457) {$a=1,\,n=1$};                % L q13_init   核字框并集 x[554.256,610.282] y[299.531,310.205]
  \node at (582.401,263.493) {$|a-1.414|\geqslant0.005$?}; % L q13_test   核字框并集 x[536.915,627.623] y[257.504,270.846]
  \node at (591.066,237.873) {是};                         % L q13_yes    核字框中心 (591.607,238.157)
  \node at (581.595,215.961) {$a=1+\dfrac{1}{1+a}$};       % L q13_assign 核字框并集 x[550.921,612.950] y[203.468,226.817]
  \node at (582.184,179.721) {$n=n+1$};                    % L q13_incr   核字框并集 x[560.259,604.946] y[174.115,183.455]
  \node at (656.393,271.773) {否};                         % L q13_no     核字框中心 (656.970,272.180)
  \node at (668.236,216.090) {输出 $n$};                   % L q13_out    核字框并集 x[652.968,683.649] y[210.139,221.480]
  \node at (668.037,184.434) {结束};                       % L q13_end    核字框中心 (668.976,184.789)
\end{tikzpicture}
